% !PRINTED none
% Chapter-opening page: no running head/page number is printed at the top of this
% page; instead the page carries the Abschnitt title block, then the §.115 heading
% and the start of the section text. (Weber suppresses the folio number on
% Abschnitt-opening pages.)

\begin{center}
\gesperrt{Fourteenth Chapter.}

\gesperrt{Groups of Linear Ternary Substitutions.}

\rule{0.3\linewidth}{0.4pt}

\S.\,115.

\gesperrt{Ternary linear substitution group of degree 168.}
\end{center}

The question of the totality of all possible finite groups of
linear substitutions in several dimensions, in particular those of
three and four dimensions, has been treated by \gesperrt{C. Jordan},
who, much as was done for the binary substitutions, has set up a
finite number of types of such groups\footnote{\gesperrt{C. Jordan}, Mémoire sur les équations différentielles linéaires à
intégrale algébrique. Crelle's Journal für Mathematik, Bd. 84 (1878). Sur
la détermination des groupes d'ordre fini etc. Atti della R. Accademia di
Napoli, Vol. VIII (1879). \gesperrt{Valentiner}, Kj\"ob. Skrift (6) V (1889).}. These general investigations
we cannot pursue here; we content ourselves with considering one
example of a ternary group in more detail.

In the chapter on the congruence groups (\S.\,66) we became
acquainted with a sequence of \gesperrt{simple} groups, of which the first,
of degree 60, proved to be isomorphic with the icosahedral group.
We now wish to consider the next of these simple groups, of
degree 168, which we have denoted by $L_7$, and attempt to
construct a ternary linear substitution group isomorphic with it.

To this we are led by the Theorem \S.\,72, I., if we seek out
four elements $\tau, \chi, \omega, \Theta$ which, under composition, satisfy the
conditions of this theorem.

We will first assume a substitution $\tau$ in the normal form
\weq{1}{\tau = \begin{pmatrix} \varepsilon_1, & 0, & 0 \\ 0, & \varepsilon_2, & 0 \\ 0, & 0, & \varepsilon_3 \end{pmatrix},}
%CONT>>
